\partie{PREMIÈRE PARTIE}{PRÉSENTATION DE L'ENTREPRISE}

% Contenu exigé par le cours (module 2, §8, premier temps) : présentation, organisation
% (organigramme), historique, secteurs, marchés, clients/fournisseurs, concurrence,
% chiffres clés, perspectives (FFOM), mutations technologiques, volonté des dirigeants.

\section{Présentation générale et historique}

\InfoEntreprise{} (\InfoEntrepriseSigle) est une société à responsabilité limitée de droit
camerounais, créée en \InfoEntrepriseCreation{} et dirigée par \InfoEntrepriseDirigeant. Elle
est spécialisée dans les \textbf{lots techniques du bâtiment} : climatisation et ventilation,
plomberie sanitaire, électricité, détection et protection incendie.

Partie d'une petite équipe de techniciens, l'entreprise s'est progressivement structurée
autour de divisions spécialisées pour répondre à des marchés de plus en plus importants :
résidences, hôtels, immeubles de bureaux et bâtiments publics. \AC{Étapes clés de l'historique :
premiers marchés, extension des activités, certifications, grands chantiers réalisés.}

Sa vocation est d'apporter à ses clients une prestation complète, de l'étude à la mise en
service, puis à la maintenance des installations.

\section{Situation géographique}

Le siège de \InfoEntrepriseSigle{} est situé à \InfoEntrepriseVille, au quartier
\InfoEntrepriseQuartier. Sa position dans un quartier accessible facilite l'approvisionnement
en matériel et le déplacement des équipes vers les chantiers. La figure~\ref{fig:localisation}
présente la localisation du siège.

\figureimg{img/fig_localisation}{0.75\textwidth}{plan de localisation du siège de l'entreprise}{fig:localisation}

\section{Missions et domaines d'activité}

\subsection{Missions}
L'entreprise a pour mission de concevoir, réaliser et entretenir les installations techniques
des bâtiments, dans le respect des règles de l'art, des délais et des budgets de ses clients.
Concrètement, elle :
\begin{itemize}
  \item réalise les études techniques et les plans d'exécution des installations ;
  \item fournit et pose les équipements et les réseaux ;
  \item effectue les essais, la mise en service et la réception des ouvrages ;
  \item assure la maintenance préventive et corrective dans le cadre de contrats d'entretien.
\end{itemize}

\subsection{Secteurs d'activité}
Le tableau~\ref{tab:activites} résume les principaux domaines d'intervention.

\begin{table}[H]\centering
\caption{domaines d'activité de l'entreprise}\label{tab:activites}
\begin{tabularx}{\textwidth}{>{\bfseries}p{3.6cm} Y}
\toprule
\entete Domaine & Prestations \\
\midrule
Climatisation & Installation de climatiseurs individuels (split, cassette), systèmes à débit de réfrigérant variable, centrales de traitement d'air, ventilation \\
Froid & Chambres froides positives et négatives, vitrines réfrigérées, maintenance frigorifique \\
Plomberie sanitaire & Réseaux d'eau froide et d'eau chaude, évacuations, appareils sanitaires, surpresseurs, production d'eau chaude \\
Électricité & Courants forts et courants faibles, éclairage, tableaux électriques \\
Sécurité incendie & Détection, extinction, robinets d'incendie armés \\
\bottomrule
\end{tabularx}
\end{table}

\subsection{Savoir-faire et atouts}
L'entreprise s'appuie sur une \textbf{offre multi-lots}, qui permet au client de traiter avec
un interlocuteur unique pour plusieurs corps d'état. Elle dispose d'un personnel
pluridisciplinaire, d'un parc d'outillage spécialisé (polyfuseurs, groupes de soudure, pompes
d'épreuve, pompes à vide) et d'une expérience des chantiers de grande hauteur.

\section{Organisation et organigramme}

\subsection{Organisation générale}
L'entreprise est organisée de façon hiérarchique. La direction générale définit la stratégie
et suit les marchés. La direction technique coordonne les études et l'exécution. Le conducteur
des travaux, également chargé d'études, organise les chantiers. Les chefs de division
(climatisation, plomberie, électricité, réseaux et gaines, génie civil) encadrent les chefs
d'équipe, eux-mêmes responsables des techniciens et des manœuvres. Un magasinier gère les
stocks et un service logistique assure les transports. L'effectif est de
\InfoEntrepriseEffectif{} personnes.

\subsection{Organigramme}
La figure~\ref{fig:organigramme} présente l'organigramme fonctionnel de l'entreprise. Notre
position de stagiaire y est indiquée en vert : nous étions rattaché à l'équipe plomberie, sous
l'autorité du chef de division plomberie.

\begin{figure}[H]\centering
% Organigramme fonctionnel (à adapter à l'entreprise réelle)
\begin{tikzpicture}[
  font=\scriptsize\bfseries, node distance=6mm and 3mm,
  bloc/.style={draw=vertfonce, fill=vertclair, rounded corners=3pt, align=center,
               minimum height=8mm, text width=2.15cm, inner sep=2pt},
  lien/.style={-{Stealth[length=2mm]}, vertfonce, line width=0.6pt}]
  \node[bloc, fill=vertfonce, text=white, text width=3.2cm] (dg) {Directeur général};
  \node[bloc, left=1.4cm of dg] (dt) {Directeur technique};
  \node[bloc, right=1.4cm of dg, text width=2.6cm] (ct) {Conducteur des travaux\\chargé d'études};
  \node[bloc, below=17mm of dg, xshift=-5.0cm] (d1) {Division\\climatisation};
  \node[bloc, right=of d1] (d2) {Division\\plomberie};
  \node[bloc, right=of d2] (d3) {Division\\électricité};
  \node[bloc, right=of d3] (d4) {Division réseaux\\et gaines};
  \node[bloc, right=of d4] (d5) {Division\\génie civil};
  \node[bloc, below=of d1] (e1) {Chef d'équipe\\climatisation};
  \node[bloc, below=of d2] (e2) {Chef d'équipe\\plomberie};
  \node[bloc, below=of d3] (e3) {Chef d'équipe\\électricité};
  \node[bloc, below=of d4] (e4) {Magasinier};
  \node[bloc, below=of d5] (e5) {Logistique\\chauffeur};
  \node[bloc, below=of e2, text width=4.4cm, xshift=1.2cm] (tech) {Techniciens};
  \node[bloc, below=of tech, text width=4.4cm] (man) {Manœuvres};
  \node[bloc, right=6mm of tech, fill=vertmoyen, text=white] (stg) {STAGIAIRE};
  \draw[lien] (dg) -- (dt);
  \draw[lien] (dg) -- (ct);
  \coordinate (bus) at ($(dg.south)+(0,-11mm)$);
  \draw[vertfonce, line width=0.6pt] (dg.south) -- (bus);
  \draw[vertfonce, line width=0.6pt] (bus -| d1.north) -- (bus -| d5.north);
  \foreach \d in {d1,d2,d3,d4,d5} \draw[lien] (bus -| \d.north) -- (\d.north);
  \foreach \d/\e in {d1/e1,d2/e2,d3/e3,d4/e4,d5/e5} \draw[lien] (\d) -- (\e);
  \draw[lien] (e2.south) -- (e2.south |- tech.north);
  \draw[lien] (tech) -- (man);
  \draw[lien, dashed, vertmoyen] (e2.south east) -- (stg.north west);
\end{tikzpicture}

\caption{organigramme fonctionnel de l'entreprise}\label{fig:organigramme}
\end{figure}

\subsection{Personnel d'encadrement rencontré}
\begin{table}[H]\centering
\caption{principaux responsables et leurs postes}\label{tab:membres}
\begin{tabularx}{0.9\textwidth}{Y Y}
\toprule
\entete Nom & Poste \\
\midrule
\InfoEntrepriseDirigeant & Directeur général \\
\AC{Nom} & Directeur technique \\
\InfoTuteurPro & \InfoTuteurProPoste \\
\AC{Nom} & Chef de division plomberie \\
\AC{Nom} & Chef de division climatisation \\
\AC{Nom} & Chef d'équipe plomberie \\
\AC{Nom} & Magasinier \\
\bottomrule
\end{tabularx}
\end{table}

\section{Environnement économique}

\subsection{Marchés et clients}
L'entreprise intervient principalement sur le \textbf{marché du bâtiment neuf} (immeubles
d'habitation, hôtels, bureaux) et sur celui de la \textbf{maintenance}. Ses clients sont des
promoteurs immobiliers privés, des maîtres d'ouvrage publics, des entreprises générales de
bâtiment qui lui sous-traitent les lots techniques, et des établissements (hôtels, banques,
cliniques) liés par des contrats d'entretien. \AC{Exemples de clients et de chantiers de
référence.}

\subsection{Fournisseurs}
Les tubes et raccords (PPR, PVC, PE, cuivre), les appareils sanitaires, les climatiseurs et les
fluides frigorigènes sont achetés auprès de distributeurs locaux, dont une partie importe le
matériel. La disponibilité et le délai de livraison de certaines références constituent un
enjeu pour le respect des plannings. \AC{Principaux fournisseurs.}

\subsection{Situation concurrentielle}
Le secteur des lots techniques est concurrentiel. Il comprend de grandes entreprises
générales, des sociétés spécialisées et de nombreux artisans informels qui pratiquent des prix
bas. \InfoEntrepriseSigle{} se différencie par son offre multi-lots, son expérience des grands
chantiers et sa capacité à assurer la maintenance après livraison.

\subsection{Chiffres clés}
\begin{table}[H]\centering
\caption{chiffres clés de l'entreprise}\label{tab:chiffres}
\begin{tabularx}{0.85\textwidth}{>{\bfseries}p{6cm} Y}
\toprule
\entete Indicateur & Valeur \\
\midrule
Forme juridique & SARL \\
Année de création & \InfoEntrepriseCreation \\
Capital social & \InfoEntrepriseCapital \\
Effectif & \InfoEntrepriseEffectif \\
Chiffre d'affaires annuel & \AC{montant} \\
Zone d'intervention & \AC{villes, régions, pays} \\
Contacts & \InfoEntrepriseContact \\
\bottomrule
\end{tabularx}
\end{table}

\section{Perspectives d'évolution}

\subsection{Analyse FFOM}
Le tableau~\ref{tab:ffom} résume notre analyse des forces, faiblesses, opportunités et menaces
de l'entreprise, établie à partir de nos observations sur le chantier et de nos échanges avec
le personnel.

\begin{table}[H]\centering
\caption{analyse FFOM de l'entreprise}\label{tab:ffom}
\begin{tabularx}{\textwidth}{Y Y}
\toprule
\entete \textbf{Forces} & \textbf{Faiblesses} \\
\midrule
Offre multi-lots (froid, plomberie, électricité) ; personnel expérimenté ; références de grands chantiers ; capacité de maintenance &
Gestion des stocks et des approvisionnements perfectible ; sécurité sur chantier insuffisamment formalisée ; dépendance à quelques gros marchés \\
\midrule
\entete \textbf{Opportunités} & \textbf{Menaces} \\
\midrule
Croissance de la construction urbaine ; demande en efficacité énergétique et en climatisation performante ; contrats de maintenance récurrents &
Concurrence des prestataires informels ; hausse du coût des matériaux importés ; retards de paiement des maîtres d'ouvrage \\
\bottomrule
\end{tabularx}
\end{table}

\subsection{Mutations technologiques et volonté des dirigeants}
Le métier évolue rapidement : généralisation des tubes de synthèse (PPR, PER, multicouche) à la
place de l'acier galvanisé, remplacement progressif des fluides frigorigènes à fort potentiel de
réchauffement global, développement des climatiseurs à technologie \emph{inverter} et des
systèmes à débit de réfrigérant variable. Les dirigeants souhaitent accompagner ces mutations en
formant leur personnel et en développant l'activité de maintenance, plus régulière que celle des
travaux neufs. \AC{Projets annoncés par la direction.}

\conclusionPartielle
\InfoEntrepriseSigle{} est une entreprise structurée, dont l'offre couvre l'ensemble des lots
techniques du bâtiment. Son organisation hiérarchique claire et la diversité de ses activités en
font un cadre propice à l'apprentissage. L'analyse de son environnement montre une entreprise
bien placée sur un marché porteur, mais qui doit renforcer sa logistique et sa culture de la
sécurité. C'est dans ce cadre que s'est déroulé notre stage, dont la deuxième partie rend compte.
